\documentclass{beamer}
\usetheme{metropolis}
\usepackage[T1]{fontenc}
\usepackage{graphicx}
\usepackage[normalem]{ulem}
\usepackage[french]{babel}
\usepackage{marvosym}
\usepackage{qrcode}
% Horizontal lists
% \usepackage[inline]{enumitem}
\usepackage{minted}

% -- ChatGPTeries pour slide dorée --%
\usepackage{tikz}
\usepackage{xcolor}

\definecolor{gold}{RGB}{201, 164, 76}
\definecolor{darkgold}{RGB}{142, 111, 35}
\definecolor{luxuryblack}{RGB}{18, 18, 18}
\definecolor{softwhite}{RGB}{242, 240, 234}
% ---- %

\usepackage{tikz}
\usepackage{xcolor}
\usetikzlibrary{positioning, calc, shapes.geometric}

% -- Icones pour la slide "langage de haut niveau" -- %
% Baguette magique : image demandee par l'utilisateur (istockphoto, a
% verifier niveau licence avant toute diffusion publique des slides).
\newcommand{\wandicon}{.}
\newcommand{\screenicon}{%
	\begin{tikzpicture}[baseline=-0.3ex, scale=0.35]
		\draw[thick, fill=black!5] (-1,0) rectangle (1,0.7);
		\node[font=\tiny] at (0,0.35) {Salut};
		\draw[thick] (-0.15,-0.35) rectangle (0.15,0);
		\draw[thick] (-0.5,-0.5) -- (0.5,-0.5);
	\end{tikzpicture}%
}
% ---- %

% Slides année précédente
\usetikzlibrary{positioning, calc, decorations.pathreplacing}

\definecolor{pillfill}{RGB}{173, 216, 240}
\definecolor{pillborder}{RGB}{70, 130, 180}
\tikzset{
	pill/.style={draw=pillborder, thick, rounded corners=8pt, fill=pillfill,
		minimum height=0.8cm, text width=3.6cm, align=center, font=\bfseries\small},
	pillsmall/.style={draw=pillborder, thick, rounded corners=6pt, fill=pillfill,
		minimum height=0.7cm, text width=2.6cm, align=center, font=\bfseries\footnotesize},
}

% -- Coloration syntaxique du code Python (compatible pdflatex) -- %
\usepackage{listings}
% Quickfix to get Listings to work properly
% From https://tex.stackexchange.com/questions/451532/recent-issues-with-lstlinebgrd-package-with-listings-after-the-latters-updates
\makeatletter
\let\old@lstKV@SwitchCases\lstKV@SwitchCases
\def\lstKV@SwitchCases#1#2#3{}
\makeatother
\usepackage{lstlinebgrd}
\makeatletter
\let\lstKV@SwitchCases\old@lstKV@SwitchCases

\lst@Key{numbers}{none}{%
	\def\lst@PlaceNumber{\lst@linebgrd}%
	\lstKV@SwitchCases{#1}%
	{none:\\%
		left:\def\lst@PlaceNumber{\llap{\normalfont
				\lst@numberstyle{\thelstnumber}\kern\lst@numbersep}\lst@linebgrd}\\%
		right:\def\lst@PlaceNumber{\rlap{\normalfont
				\kern\linewidth \kern\lst@numbersep
				\lst@numberstyle{\thelstnumber}}\lst@linebgrd}%
	}{\PackageError{Listings}{Numbers #1 unknown}\@ehc}}
\makeatother
\definecolor{codekw}{RGB}{0, 0, 180}
\definecolor{codestr}{RGB}{20, 120, 40}
\definecolor{codecom}{RGB}{120, 120, 120}
\lstdefinestyle{python}{
	language=Python,
	basicstyle=\ttfamily\footnotesize,
	keywordstyle=\color{codekw}\bfseries,
	stringstyle=\color{codestr},
	commentstyle=\color{codecom}\itshape,
	showstringspaces=false,
	keepspaces=true,
	columns=fullflexible,
	aboveskip=0pt,
	belowskip=0pt,
	literate={é}{{\'e}}1 {è}{{\`e}}1 {ê}{{\^e}}1 {à}{{\`a}}1 {ç}{{\c c}}1
	{É}{{\'E}}1 {'}{{\textquotesingle}}1,
}
\lstdefinestyle{pythonsmall}{style=python,basicstyle=\ttfamily\scriptsize}

% Convention du cours pour une reponse de Python (a la suite d'un print, etc.) :
% prefixe ">" + encadre sur fond clair.
\lstdefinestyle{pythonoutput}{
	style=python,
	basicstyle=\ttfamily\footnotesize\color{black!70},
	backgroundcolor=\color{black!6},
	frame=single,
	rulecolor=\color{black!35},
	framerule=0.4pt,
	framexleftmargin=5pt,
	framexrightmargin=5pt,
	framextopmargin=3pt,
	framexbottommargin=3pt,
	xleftmargin=5pt,
	xrightmargin=5pt,
	aboveskip=3pt,
	belowskip=3pt,
}

% Un morceau de code annote : #1 plage de lignes, #2 1 si c'est la partie
% commentee (fond bleu, couleurs normales) sinon 0 (grise), #3 aboveskip.
\newcommand{\annoline}[3]{%
	\ifnum#2=1
	\lstinputlisting[style=pythonsmall, backgroundcolor=\color{blue!15}, #3, linerange={#1}]{code/premier_prog_2.py}%
	\else
	\lstinputlisting[style=pythonsmall, basicstyle=\ttfamily\scriptsize\color{black!45}, keywordstyle=\color{black!45}, stringstyle=\color{black!45}, #3, linerange={#1}]{code/premier_prog_2.py}%
	\fi
}

% Version generique de \annoline : #1 fichier, #2 plage de lignes,
% #3 1 si c'est la ligne commentee (fond bleu) sinon 0 (grisee).
\newcommand{\valline}[3]{%
	\ifnum#3=1
	\lstinputlisting[style=pythonsmall, backgroundcolor=\color{blue!15}, linerange={#2}]{#1}%
	\else
	\lstinputlisting[style=pythonsmall, basicstyle=\ttfamily\scriptsize\color{black!45}, keywordstyle=\color{black!45}, stringstyle=\color{black!45}, linerange={#2}]{#1}%
	\fi
}

% Slide "focus sur une construction" : etape 1 = programme entier, la partie
% commentee sur fond bleu (texte noir), le reste gris ; etape 2 = le reste a
% disparu, la partie commentee reste sur fond bleu (texte noir).
\newcommand{\focusframe}[5]{%
	\begin{frame}[t]{#1}
		\only<1>{\lstinputlisting[
			style=pythonsmall,
			basicstyle=\ttfamily\scriptsize\color{black!45},
			keywordstyle=\color{black!45}, stringstyle=\color{black!45},
			emph={#2}, emphstyle={\color{black}\bfseries},
			linebackgroundcolor={\color{white}#3},
			]{code/premier_prog_2.py}}%
		\only<2>{\lstinputlisting[
			style=pythonsmall,
			emph={#2}, emphstyle={\color{black}\bfseries},
			linebackgroundcolor={\color{white}#3},
			]{code/#4}}%
		\vfill
		\centering\small #5
	\end{frame}%
}

\title{Algorithmique et Programmation 1}
\author{A. Boussidan, L. Exibard, P. Yvonnet}
\date{28 Septembre 2026}

\begin{document}
	
	\begin{frame}
		\maketitle
	\end{frame}
	
	\begin{frame}{Un retour sur les boucles while}
		
		Que faire si on veut répéter une instruction si une condition reste vraie ?
		
		\begin{columns}
			\begin{column}{0.7\textwidth}
				\scalebox{0.8}{%
					\begin{tikzpicture}[
						every node/.style={font=\small}	,
						process/.style={
							rectangle,
							draw,
							rounded corners,
							minimum width=3cm,
							minimum height=1cm,
							align=center
						},
						decision/.style={
							diamond,
							draw,
							aspect=2,
							minimum width=3cm,
							minimum height=1.2cm,
							align=center
						},
						arrow/.style={
							->,
							>=stealth,
							thick
						}
						]
						
						% Blocs
						\node[process] (debut) {Début};
						
						\node[decision, below=1.2cm of debut] (condition) {Condition ?};
						
						\node[process, below=1.5cm of condition] (bloc) {Bloc};
						
						\node[process, below=1.5cm of bloc] (suite) {Suite du programme};
						
						
						% Début -> Condition
						\draw[arrow]
						(debut) -- (condition);
						
						
						% Condition -> Bloc : OUI
						\draw[arrow]
						(condition.south)
						-- node[right] {oui}
						(bloc.north);
						
						
						% Bloc -> Condition : retour de la boucle
						\draw[arrow]
						(bloc.west)
						-- ++(-2cm,0)
						|- node[left, pos=0.25] {}
						(condition.west);
						
						
						% Condition -> Suite : NON
						\draw[arrow]
						(condition.east)
						to[out=0, in=0]
						node[right] {non}
						(suite.east);
						
					\end{tikzpicture}
					
				
				}
			\end{column}
			\begin{column}{0.3\textwidth}
				\lstinputlisting[style=pythonsmall]{code/while.py}%
			\end{column}
		\end{columns}
	\end{frame}
	
	\begin{frame}{Un peu de vocabulaire}
		\begin{block}{Itération}
			\vspace{0pt}
			Une exécution complète du corps de la boucle
		\end{block}
		
		\pause 
		\begin{block}{Variable ou indice de boucle}
			\vspace{0pt}
			La variable qui compte le nombre de tours de boucles ; dans une boucle while simple, c'est sur elle que porte la condition.\\
			Souvent notée \mintinline{python}{i}
			ou 
			\mintinline{python}{j}
		\end{block}
		\pause 
		\begin{block}{Condition de continuation}
			La condition qui est vérifiée à chaque début de boucle : si elle est vraie, on continue. 
		\end{block}
		\pause
		\begin{block}{Condition d'arrêt}
			L'inverse de la condition de continuation
		\end{block}
	\end{frame}
	
	\begin{frame}[fragile]{Sur un exemple}
		\lstinputlisting[style=python]{code/while_annote.py}
	\end{frame}
	
	\begin{frame}[fragile]{Un exemple familier...}
				\lstinputlisting[style=pythonsmall]{code/lievre_tortue_CM.py}
	\pause
	\begin{center}
		\textbf{M'sieur, ça marche pas !}
	\end{center}
	\end{frame}
	\begin{frame}[fragile]{Un exemple familier...corrigé}
		\lstinputlisting[style=pythonsmall]{code/lievre_tortue_CM_fixed.py}
		\begin{center}
			\textbf{Attention à la boucle infinie}
		\end{center}
	\end{frame}
	\begin{frame}{Comment débugguer une boucle while}
		\begin{itemize}
			\item Ajouter des prints dans son code
			\pause\item Exécuter son code pas à pas ( à la main, sur PythonTutor, avec le débbugueur)
			\pause\item S'il y a une boucle infinie : se concentrer sur les variables du test de continuation 
				\begin{itemize}
					\item Sont-elles bien initialisées ? 
					\item Quand sont-elles mises à jour ? 
					\item Pourquoi est-ce que la boucle finit par s'arrêter ? 
				\end{itemize}
		\end{itemize}
	\end{frame}
	\begin{frame}{Boucles while - Cas d'utilisation }
		\begin{block}{Rappel - être un chef !}
			\vspace{0 pt}
			On commence toujours par formuler l'algorithme.
		\end{block}
		
	\end{frame}
\end{document}